% !TEX root = ../main.tex
\chapter*{LỜI CẢM ƠN}
\addcontentsline{toc}{chapter}{LỜI CẢM ƠN}

Trong quá trình học tập và thực hiện đề tài \textbf{``Nghiên cứu, thiết kế bộ điều khiển cân bằng quadrotor''}, chúng em có cơ hội tổng hợp kiến thức đã học và vận dụng vào một bài toán kỹ thuật cụ thể. Việc nghiên cứu nguyên lý hoạt động của quadrotor, cảm biến và bộ điều khiển góp phần làm rõ mối liên hệ giữa lý thuyết với thực tế. Bên cạnh nỗ lực của các thành viên, sự hướng dẫn của thầy cô, sự quan tâm của gia đình và sự hỗ trợ của bạn bè có ý nghĩa quan trọng đối với quá trình thực hiện đề tài.

Trước hết, chúng em xin bày tỏ lòng biết ơn sâu sắc tới \textbf{\SupervisorName}, giảng viên hướng dẫn đề tài. Sự hướng dẫn của thầy là cơ sở để chúng em xác định mục tiêu, lựa chọn nội dung nghiên cứu và tổ chức công việc phù hợp với yêu cầu của đồ án. Những chỉ dẫn của thầy về cách phân tích vấn đề, sử dụng tài liệu và xây dựng phương án thiết kế có ý nghĩa thiết thực đối với việc nghiên cứu bài toán cân bằng quadrotor. Chúng em trân trọng những góp ý ấy cùng các bài học về sự cẩn thận trong tính toán và cách trình bày nội dung nghiên cứu rõ ràng, có căn cứ.

Chúng em xin trân trọng cảm ơn quý thầy cô Khoa Cơ điện tử, Trường Cơ khí -- Ô tô, Đại học Công nghiệp Hà Nội đã giảng dạy và truyền đạt kiến thức trong suốt quá trình học tập. Các kiến thức về cơ khí, điện tử, đo lường, điều khiển và lập trình là nền tảng để chúng em tiếp cận và thực hiện đề tài. Phương pháp học tập, phân tích và giải quyết vấn đề được rèn luyện qua các học phần cũng góp phần hình thành cách làm việc chủ động, có hệ thống. Chúng em biết ơn công sức giảng dạy của thầy cô cùng những bài học về tinh thần tự học, sự nghiêm túc và trách nhiệm đối với công việc.

Chúng em xin gửi lời cảm ơn tới gia đình vì sự quan tâm, động viên và tin tưởng trong quá trình học tập. Sự đồng hành của gia đình là nguồn khích lệ để chúng em kiên trì nghiên cứu, khắc phục khó khăn và cố gắng thực hiện nhiệm vụ của mình. Chúng em trân trọng những tình cảm ấy và xem đó là động lực để tiếp tục học tập, rèn luyện, nâng cao ý thức trách nhiệm với bản thân và công việc.

Bên cạnh đó, chúng em xin cảm ơn bạn bè đã chia sẻ và hỗ trợ trong thời gian học tập, thực hiện đề tài. Những trao đổi trong học tập mang đến các góc nhìn khác nhau về vấn đề nghiên cứu, đồng thời giúp chúng em nhận ra những nội dung cần bổ sung và làm rõ. Chúng em trân trọng sự quan tâm và tinh thần giúp đỡ lẫn nhau trong suốt quá trình học tập.

Do kiến thức và kinh nghiệm thực tế còn hạn chế, báo cáo có thể còn thiếu sót trong việc phân tích, lựa chọn phương án và trình bày nội dung. Chúng em kính mong nhận được những nhận xét, góp ý của quý thầy cô để tiếp tục hoàn thiện báo cáo và củng cố kiến thức chuyên môn. Những ý kiến đóng góp ấy cũng là cơ sở để chúng em xác định các vấn đề cần nghiên cứu thêm và nâng cao khả năng vận dụng kiến thức vào thực tế.

Chúng em xin chân thành cảm ơn.
